% ============================================================================
% SPANISCH LANGENTWURF - DS 3: Bucharbeit + Assessment-Ankündigung
% ============================================================================
% Basiert auf: /templates/spanisch/langentwurf-spanisch.tex
% Fachfarbe: #c41e3a (Karminrot)
%
% KURS: Jahrgangsübergreifend Spanisch Spätbeginner (Kl. 10, 11, 12)
% SCHÜLERZAHL: 10 Schüler
% LEHRWERK: A_tope.com (Cornelsen)
%
% PFLICHT: Tier 1 (vollständig)
% ============================================================================

\documentclass[11pt,a4paper]{article}

% ============================================================================
% PACKAGES
% ============================================================================
\usepackage[utf8]{inputenc}
\usepackage[T1]{fontenc}
\usepackage[ngerman]{babel}
\usepackage{geometry}
\usepackage{fancyhdr}
\usepackage{lastpage}
\usepackage[table]{xcolor}
\usepackage{tcolorbox}
\usepackage{enumitem}
\usepackage{tabularx}
\usepackage{longtable}
\usepackage{array}
\usepackage{booktabs}
\usepackage{hyperref}
\usepackage{ifthen}
\usepackage{amssymb}

% ============================================================================
% KONFIGURATION
% ============================================================================

% --- TIER (immer 1 für Spanisch) ---
\newcommand{\plantier}{1}

% --- FACH ---
\newcommand{\fachid}{spanisch}
\newcommand{\fach}{Spanisch}
\newcommand{\fachfarbe}{c41e3a}

% ============================================================================
% VARIABLEN
% ============================================================================

% --- Metadaten ---
\newcommand{\jahrgangsstufe}{10/11/12 (Spätbeginner)}
\newcommand{\schulart}{Gesamtschule}
\newcommand{\stundenthema}{Bucharbeit + Assessment-Ankündigung}
\newcommand{\reihenthema}{Unterrichtsvorbereitung Beobachtung Feb. 2026}
\newcommand{\stundennummer}{3}
\newcommand{\reihenstunden}{10}
\newcommand{\unterrichtsdatum}{12.01.2026}
\newcommand{\unterrichtszeit}{Doppelstunde (90 Min.)}
\newcommand{\raum}{---}

% --- Lehrkraft ---
\newcommand{\lehrkraft}{N.N.}
\newcommand{\ausbildungsstand}{Lehrkraft}

% --- Lerngruppe ---
\newcommand{\lerngruppengroesse}{10}
\newcommand{\maennlich}{---}
\newcommand{\weiblich}{---}
\newcommand{\divers}{---}

% --- Kompetenzen/Niveau ---
\newcommand{\gerniveau}{A1--B1 (differenziert)}
\newcommand{\lehrwerk}{A\_tope.com (Cornelsen)}
\newcommand{\lehrwerkseite}{S. 32--33 (A) / S. 105--106 (B)}

% ============================================================================
% LAYOUT
% ============================================================================
\geometry{left=2.5cm, right=2.5cm, top=2.5cm, bottom=2.5cm}
\definecolor{fach}{HTML}{\fachfarbe}
\definecolor{gray}{HTML}{666666}
\definecolor{lightgray}{HTML}{f5f5f5}
\definecolor{successgreen}{HTML}{2a7a4b}
\definecolor{warningorange}{HTML}{b35c00}

\tcbuselibrary{skins,breakable}

% Farbige Boxen
\newtcolorbox{infobox}[1][]{
    colback=lightgray,
    colframe=fach,
    fonttitle=\bfseries,
    title=#1,
    breakable
}

\newtcolorbox{scorebox}{
    colback=white,
    colframe=fach,
    fonttitle=\bfseries,
    title=Didaktik-Score,
    breakable
}

\newtcolorbox{warnbox}{
    colback=white,
    colframe=warningorange,
    fonttitle=\bfseries,
    title=WICHTIG: Assessment-Ankündigung,
    breakable
}

% Kopf-/Fußzeile
\pagestyle{fancy}
\fancyhf{}
\fancyhead[L]{\small\textcolor{gray}{\fach{} -- Jg. \jahrgangsstufe{} -- \stundenthema}}
\fancyhead[R]{\small\textcolor{gray}{Langentwurf Tier 1}}
\fancyfoot[C]{\small\thepage{} / \pageref{LastPage}}
\renewcommand{\headrulewidth}{0.4pt}
\renewcommand{\footrulewidth}{0pt}

% ============================================================================
% DOKUMENT
% ============================================================================
\begin{document}

% ============================================================================
% DECKBLATT
% ============================================================================
\begin{titlepage}
    \centering
    \vspace*{2cm}

    {\large\textcolor{gray}{\schulart}}\\[2cm]

    {\Huge\textcolor{fach}{\textbf{Langentwurf}}}\\[0.5cm]
    {\large Tier 1 -- Vollständig}\\[2cm]

    {\LARGE\textbf{\stundenthema}}\\[0.5cm]
    {\large im Rahmen der Unterrichtsreihe}\\[0.3cm]
    {\Large\textit{\reihenthema}}\\[2cm]

    \begin{tabular}{rl}
        \textbf{Fach:} & \fach{} (\gerniveau) \\[0.3cm]
        \textbf{Klasse:} & \jahrgangsstufe \\[0.3cm]
        \textbf{Lehrwerk:} & \lehrwerk \\[0.3cm]
        \textbf{Seiten:} & \lehrwerkseite \\[0.3cm]
        \textbf{Datum:} & Montag, \unterrichtsdatum \\[0.3cm]
        \textbf{Zeit:} & \unterrichtszeit{} (Doppelstunde) \\[0.3cm]
    \end{tabular}

    \vfill

    {\large Lehrkraft: \lehrkraft}\\[0.3cm]
    {\normalsize\textcolor{gray}{\ausbildungsstand}}\\[1cm]

\end{titlepage}

% ============================================================================
% INHALTSVERZEICHNIS
% ============================================================================
\tableofcontents
\newpage

% ============================================================================
% 1. BEDINGUNGSANALYSE
% ============================================================================
\section{Bedingungsanalyse}

\subsection{Lerngruppe}
\begin{tabular}{ll}
    Kursgröße: & \lerngruppengroesse{} SuS (jahrgangsübergreifend) \\
    Jahrgangsstufen: & Klasse 10 (5 SuS), Klasse 11 (2 SuS), Klasse 12 (3 SuS) \\
    Sprachniveau: & A1--B1 (stark differenziert) \\
    Fremdsprache: & Spätbeginner (Spanisch ab Kl. 10/11) \\
    Gruppierung: & Gruppe A (Kl. 10 + L.K.), Gruppe B (Kl. 12 + stärkere Kl. 11) \\
\end{tabular}

\subsection{Differenzierung (Lernjahr)}

\textbf{HINWEIS:} Aufgrund der jahrgangsübergreifenden Struktur erfolgt Differenzierung primär durch Gruppenarbeit (A vs. B), nicht durch Niveaustufen auf dem Material.

\begin{tabular}{|l|p{5cm}|p{5cm}|}
\hline
& \textbf{Gruppe A (LJ 1)} & \textbf{Gruppe B (LJ 2--3)} \\
\hline
\textbf{Niveau} & A1--A2 & B1 \\
\hline
\textbf{Inhalt} & Basisgrammatik (ir, poder, querer, tener que) & Komplexe Strukturen (Argumentation, Konjunktionen) \\
\hline
\textbf{Seiten} & S. 32--33 & S. 105--106 \\
\hline
\textbf{Hilfen} & Konjugationstabellen, Wortschatzliste & Wörterbuch, eigenständige Arbeit \\
\hline
\end{tabular}

\subsection{Schülerprofile (Gruppe A)}

\begin{tabular}{|l|c|p{3cm}|p{3cm}|p{4cm}|}
\hline
\textbf{Kürzel} & \textbf{Niv.} & \textbf{Stärken} & \textbf{Schwächen} & \textbf{Strategie} \\
\hline
E.H. & 15 NP & Franz. Hintergrund, immer aktiv & --- & Peer-Tutor, anspruchsvollere Aufgaben \\
\hline
L.N. & $\sim$13 NP & Proaktiv, versteht schnell & Etwas faul, ablenkbar & In Bewegung halten, Verantwortung \\
\hline
C.P. & 12 NP & Gute Aussprache, ruhig & Braucht Zeit & Klare Anweisungen, Wartezeit \\
\hline
V.S. & $\sim$11 NP & Konstant, fleißig & Wenig Selbstvertrauen & Kleine Erfolge, Lob \\
\hline
L.K. & 10 NP & Versteht gut (Kl. 11) & Ablenkbar, weniger fleißig & Struktur, klare Erwartungen \\
\hline
H.P. & 10 NP & Bemüht & Nachzügler, Grundlagenlücken & Scaffolding, Partner E.H. \\
\hline
A.A. & $\sim$10 NP & Fleißig, gute Tests & Unsicher, schaut zu E.H. & Ermutigung, getrennt von E.H. \\
\hline
A.S. & 9 NP & Macht Fortschritte & Unsicher, fragend & Positive Verstärkung \\
\hline
\end{tabular}

\subsection{Schülerprofile (Gruppe B)}

Gruppe B umfasst 2--3 Schüler (Kl. 12 + stärkere Kl. 11), die auf B1-Niveau arbeiten. Details im Beobachtungsplan.

% ============================================================================
% 2. SACHANALYSE
% ============================================================================
\section{Sachanalyse}

\subsection{Sprachliche Strukturen -- Gruppe A}

\subsubsection{Verb \textit{ir} (gehen/fahren)}
\begin{tabular}{ll}
    voy, vas, va, vamos, vais, van \\
\end{tabular}

\textbf{Konstruktion:} ir a + Ort (¿Adónde vais? -- Vamos al cine.)

\subsubsection{Modalverben}
\begin{itemize}
    \item \textbf{poder} (können): puedo, puedes, puede, podemos, podéis, pueden
    \item \textbf{querer} (wollen): quiero, quieres, quiere, queremos, queréis, quieren
    \item \textbf{tener que} (müssen): tengo que, tienes que, tiene que...
\end{itemize}

\subsubsection{Phonetik: [b] vs. [\ss]}
\begin{itemize}
    \item \textbf{[b]} (okklusiv): Am Wortanfang, nach \textit{m/n} (beber, vamos, enviar)
    \item \textbf{[\ss]} (frikativ): Zwischen Vokalen, nach anderen Konsonanten (saber, libro)
\end{itemize}

\subsection{Sprachliche Strukturen -- Gruppe B}

\subsubsection{Thema: Arbeiten im Ausland}
\begin{itemize}
    \item \textbf{Wortfeld:} el trabajo, la profesión, el empleo, las cualificaciones
    \item \textbf{Gründe für Auslandsaufenthalt:} ganar experiencia, aprender idiomas, conocer otras culturas
    \item \textbf{Texte:} Carla (arbeitet in Barcelona), Jan (Praktikum in Madrid)
\end{itemize}

\subsubsection{Kompetenzfokus}
\begin{itemize}
    \item Leseverstehen: Textverständnis (dilo con otras palabras)
    \item Wörterbucharbeit: Strategien zum Nachschlagen
    \item Sprechen: Gründe formulieren (porque, ya que, para + Infinitiv)
\end{itemize}

\subsection{Kontrastive Betrachtung}

\begin{tabular}{|l|p{5cm}|p{5cm}|}
\hline
\textbf{Aspekt} & \textbf{Deutsch} & \textbf{Spanisch} \\
\hline
Modalverben + Infinitiv & Ich kann gehen & Puedo ir (Infinitiv am Ende) \\
\hline
[b]/[v] Unterscheidung & Klar unterschieden & Kein Unterschied! \\
\hline
Präposition ``nach'' & nach + Ort & a + Ort (ir a Madrid) \\
\hline
\end{tabular}

% ============================================================================
% 3. DIDAKTISCHE ANALYSE
% ============================================================================
\section{Didaktische Analyse}

\subsection{Stellung in der Sequenz}
Dies ist die \textbf{3. Doppelstunde} der Sequenz zur Unterrichtsvorbereitung Beobachtung (10 Doppelstunden gesamt).

\begin{itemize}
    \item DS 1--2: Bucharbeit (parallel)
    \item \textbf{DS 3: Bucharbeit + Assessment-Ankündigung} $\leftarrow$ Diese Stunde
    \item DS 4: Letzte Bucharbeit vor Assessment
    \item DS 5: \textbf{ASSESSMENT} (sonstige Note)
    \item DS 6--10: Brückenthema, Konvergenz, Beobachtung
\end{itemize}

\subsection{Kompetenzziele}

\subsubsection{Gruppe A}
\begin{itemize}
    \item \textbf{Grammatik:} Verb \textit{ir} konjugieren und mit Ortsangaben verwenden
    \item \textbf{Grammatik:} Modalverben \textit{poder/querer/tener que} im Kontext anwenden
    \item \textbf{Phonetik:} Unterscheidung [b]/[\ss] erkennen und anwenden
    \item \textbf{Kommunikativ:} Über Freizeitaktivitäten und Häufigkeit sprechen
\end{itemize}

\subsubsection{Gruppe B}
\begin{itemize}
    \item \textbf{Leseverstehen:} Texte über Arbeit im Ausland verstehen
    \item \textbf{Wortschatz:} Fachvokabular Beruf/Arbeit erweitern
    \item \textbf{Methodik:} Wörterbucharbeit vertiefen
    \item \textbf{Kommunikativ:} Gründe für Entscheidungen formulieren
\end{itemize}

\subsection{Legitimation}

\textbf{Rahmenplan Spanisch MV (Spätbeginner):}
\begin{itemize}
    \item \textbf{Sprechen:} Über Alltag, Freizeit, Pläne sprechen (A)
    \item \textbf{Lesen:} Authentische Texte verstehen (B)
    \item \textbf{Sprachmittlung:} Wörterbucharbeit (B)
\end{itemize}

% ============================================================================
% 4. STUNDENZIELE
% ============================================================================
\section{Stundenziele}

\subsection{Grobziel}
Die SuS erweitern ihre \textbf{grammatische Kompetenz} (Gruppe A: ir + Modalverben; Gruppe B: Leseverstehen) und werden auf das \textbf{Assessment am Mo 19.01.} vorbereitet, indem sie die Bewertungskriterien kennenlernen.

\subsection{Feinziele (operationalisiert)}

\begin{tabular}{|c|p{7cm}|c|c|}
\hline
\textbf{Nr.} & \textbf{Die SuS können...} & \textbf{AFB} & \textbf{Gruppe} \\
\hline
FZ1 & das Verb \textit{ir} in allen Personen korrekt konjugieren & I & A \\
\hline
FZ2 & \textit{ir a + Ort} verwenden, um Ziele zu benennen & II & A \\
\hline
FZ3 & Modalverben \textit{poder/querer/tener que} im Lückentext einsetzen & II & A \\
\hline
FZ4 & die Aussprache [b]/[\ss] anhand von Beispielen unterscheiden & I & A \\
\hline
FZ5 & Gründe für Auslandsaufenthalte aus dem Text entnehmen & II & B \\
\hline
FZ6 & unbekannte Wörter mit dem Wörterbuch erschließen & II & B \\
\hline
FZ7 & Textinhalte mit eigenen Worten wiedergeben & III & B \\
\hline
FZ8 & die Bewertungskriterien für das Assessment benennen & I & alle \\
\hline
\end{tabular}

% ============================================================================
% 5. METHODISCHE ANALYSE
% ============================================================================
\section{Methodische Analyse}

\subsection{Phasierung (Doppelstunde = 2 $\times$ 45 Min.)}

\subsubsection{1. Stunde (45 Min.)}
\begin{enumerate}
    \item \textbf{Einstieg (5 Min.):} Begrüßung, Tagesübersicht, kurzer Warm-up
    \item \textbf{Erarbeitung A (15 Min.):} Aufg. 4a/b -- ir-Konjugation + Zahlen
    \item \textbf{Erarbeitung B (15 Min.):} Aufg. 1--2 -- Gründe Ausland + Wörterbuch
    \item \textbf{Übung A (10 Min.):} Aufg. 5 -- Modalverben Lückentext
    \item \textbf{Puffer (5 Min.):} Lehrkraft wechselt zwischen Gruppen
\end{enumerate}

\subsubsection{2. Stunde (45 Min.)}
\begin{enumerate}
    \item \textbf{Wiederholung (5 Min.):} Kurze Kontrolle der Hausaufgaben
    \item \textbf{Phonetik A (10 Min.):} Aufg. 6 -- Aussprache [b]/[\ss]
    \item \textbf{Textarbeit B (15 Min.):} Aufg. 3a -- Textverständnis
    \item \textbf{Sicherung (5 Min.):} Ergebnisvergleich
    \item \textcolor{fach}{\textbf{Assessment-Ankündigung (10 Min.):}} Kriterien, Format, Termine
\end{enumerate}

\subsection{Medien}
\begin{itemize}
    \item Lehrwerk: A\_tope.com, S. 32--33 (A) / S. 105--106 (B)
    \item Tafel: Konjugationstabelle \textit{ir}
    \item Wörterbücher (für Gruppe B)
    \item Bewertungsbögen (Kopien für alle SuS)
    \item iPad/TV als Backup (Projektion)
\end{itemize}

% ============================================================================
% 6. VERLAUFSPLANUNG
% ============================================================================
\section{Verlaufsplanung}

\textbf{1. Stunde (45 Min.)}

\begin{longtable}{|p{0.8cm}|p{1.5cm}|p{4.5cm}|p{3cm}|p{1.3cm}|p{1.8cm}|}
\hline
\textbf{Zeit} & \textbf{Phase} & \textbf{Lehrerhandeln} & \textbf{Schülerhandeln} & \textbf{SF} & \textbf{Medien} \\
\hline
\endhead

0--5 & Einstieg & Begrüßung, Tagesplan vorstellen: ``Heute arbeiten wir an den Buchseiten weiter. Am Ende der Stunde habe ich wichtige Infos zum Assessment!'' & Zuhören, Fragen stellen & Plenum & Tafel \\
\hline

5--20 & Erarbei- tung A & \textbf{Gruppe A:} Aufg. 4a anleiten -- Zahlen + ir-Formen zuordnen. Konjugationstabelle an Tafel. Aufg. 4b: ¿Adónde vais? Frage-Antwort-Muster. & A: Tabelle ausfüllen, Formen üben, Sätze bilden & EA, PA & Buch, Tafel \\
\hline

5--20 & Erarbei- tung B & \textbf{Gruppe B:} Aufg. 1 -- Gründe für Auslandsaufenthalt sammeln. Aufg. 2 -- Wörterbucharbeit einleiten. & B: Brainstorming, Wörterbuch nutzen & EA, PA & Buch, Wörterbuch \\
\hline

20--35 & Übung A & \textbf{Gruppe A:} Aufg. 5 -- Lückentext mit poder/querer/tener que. L unterstützt bei Fragen. & A: Lückentext bearbeiten, mit Partner vergleichen & EA, PA & Buch \\
\hline

20--35 & Fort- setzung B & \textbf{Gruppe B:} Aufg. 2 fortführen, Ergebnisse sichern. L wechselt zu B. & B: Vokabeln notieren, verstehen & EA & Buch, Heft \\
\hline

35--40 & Zwischen- sicherung & Kurze Abfrage: ``Gruppe A -- Wie konjugiert man ir?'' Gruppe B arbeitet selbstständig. & A: Antworten, B: weiterarbeiten & Plenum (A) & -- \\
\hline

40--45 & Puffer & Bei Gruppe A letzte Fragen klären. Gruppe B auf Aufg. 3a vorbereiten. & Aufräumen, Orientierung & -- & -- \\
\hline
\end{longtable}

\vspace{1em}

\textbf{2. Stunde (45 Min.)}

\begin{longtable}{|p{0.8cm}|p{1.5cm}|p{4.5cm}|p{3cm}|p{1.3cm}|p{1.8cm}|}
\hline
\textbf{Zeit} & \textbf{Phase} & \textbf{Lehrerhandeln} & \textbf{Schülerhandeln} & \textbf{SF} & \textbf{Medien} \\
\hline
\endhead

0--5 & Wieder- holung & Kurze HA-Kontrolle (falls vorhanden). Rückblick 1. Stunde. & Zeigen, korrigieren & Plenum & -- \\
\hline

5--15 & Phonetik A & \textbf{Gruppe A:} Aufg. 6 -- Aussprache [b] vs. [\ss]. L demonstriert, SuS sprechen nach. Chorisches Sprechen. & A: Nachsprechen, Regeln erkennen & Plenum (A), EA & Buch, Tafel \\
\hline

5--20 & Text- arbeit B & \textbf{Gruppe B:} Aufg. 3a -- Leseverstehen ``Dilo con otras palabras''. Text Carla/Jan lesen, mit eigenen Worten zusammenfassen. & B: Lesen, Notizen machen, formulieren & EA & Buch, Heft \\
\hline

15--25 & Festigung A & Gruppe A übt Aussprache in Paaren. L wechselt zu Gruppe B. & A: Partnerübung Aussprache & PA & Buch \\
\hline

20--30 & Sicherung B & Gruppe B: Ergebnisse besprechen. ``Wer kann Carlas Gründe zusammenfassen?'' & B: Vorstellen, ergänzen & Plenum (B) & -- \\
\hline

30--35 & Sicherung A & Gruppe A: Stichprobenartige Kontrolle Lückentext. Lösungen kurz durchgehen. & A: Vergleichen, korrigieren & Plenum (A) & Buch \\
\hline

\rowcolor{lightgray}
35--45 & \textbf{Assess- ment-Info} & \textcolor{fach}{\textbf{ALLE:}} Assessment-Ankündigung (siehe Abschnitt 7). Format, Kriterien, Termin erklären. Bewertungsbögen austeilen. Fragen klären. & Zuhören, Fragen stellen, notieren & Plenum & Bewertungs- bögen \\
\hline
\end{longtable}

\textbf{Legende:} EA = Einzelarbeit, PA = Partnerarbeit, SF = Sozialform

% ============================================================================
% 7. ASSESSMENT-ANKÜNDIGUNG (DETAIL)
% ============================================================================
\section{Assessment-Ankündigung (Kerninhalt DS 3)}

\begin{warnbox}
\textbf{Termin:} Montag, 19.01.2026 (DS 5)

\textbf{Format:}
\begin{itemize}
    \item \textbf{Gruppe A:} Mündliche Partnerprüfung ``Mi barrio'' (3--4 Min. pro Paar)
    \item \textbf{Gruppe B:} Kurzvortrag ``Mi profesión ideal'' (3--4 Min. pro Person)
\end{itemize}

\textbf{In dieser Stunde:}
\begin{enumerate}
    \item Format und Ablauf erklären
    \item Bewertungskriterien vorstellen (Bögen austeilen!)
    \item Fragen klären
    \item Vorbereitungszeit bis DS 4 (Di 13.01.) betonen
\end{enumerate}
\end{warnbox}

\subsection{Ablauf der Ankündigung (10 Min.)}

\begin{tabular}{|c|p{10cm}|}
\hline
\textbf{Min.} & \textbf{Inhalt} \\
\hline
1--2 & ``Nächste Woche Montag ist Assessment -- eine sonstige Note.'' \\
\hline
2--4 & \textbf{Gruppe A:} Partnergespräch über euer Stadtviertel. Ihr fragt euch gegenseitig: ¿Hay...? ¿Cómo es...? ¿Adónde vas...? \\
\hline
4--6 & \textbf{Gruppe B:} Kurzvortrag ``Mein Traumberuf''. Ihr erklärt: Welcher Beruf? Warum? Welche Fähigkeiten? \\
\hline
6--8 & Bewertungsbögen austeilen und erklären (Kriterien durchgehen) \\
\hline
8--10 & Fragen beantworten. ``Am Dienstag (DS 4) üben wir nochmal!'' \\
\hline
\end{tabular}

\subsection{Bewertungskriterien Gruppe A}

\begin{tabular}{|l|c|p{6cm}|}
\hline
\textbf{Kriterium} & \textbf{Gewicht} & \textbf{Beschreibung} \\
\hline
Verwendung ser/estar/hay & 30\% & Korrekte Unterscheidung und Anwendung \\
\hline
Wortschatz (Stadtviertel) & 25\% & Orte, Gebäude, Beschreibungen \\
\hline
Grammatische Korrektheit & 20\% & Verbformen, Genus, Numerus \\
\hline
Flüssigkeit + Verständlichkeit & 25\% & Sprechtempo, Aussprache, Klarheit \\
\hline
\end{tabular}

\subsection{Bewertungskriterien Gruppe B}

\begin{tabular}{|l|c|p{6cm}|}
\hline
\textbf{Kriterium} & \textbf{Gewicht} & \textbf{Beschreibung} \\
\hline
Argumentation + Konjunktionen & 30\% & Begründungen mit ya que, porque, además \\
\hline
Fachvokabular (Berufe) & 25\% & Berufsspezifische Begriffe \\
\hline
Struktur + Kohärenz & 20\% & Logischer Aufbau, roter Faden \\
\hline
Flüssigkeit + Aussprache & 25\% & Freies Sprechen, klare Aussprache \\
\hline
\end{tabular}

% ============================================================================
% 8. ERWARTETE SCHWIERIGKEITEN
% ============================================================================
\section{Erwartete Schwierigkeiten}

\begin{tabular}{|p{4cm}|p{4cm}|p{5cm}|}
\hline
\textbf{Schwierigkeit} & \textbf{Betroffene SuS} & \textbf{Reaktion} \\
\hline
Verwechslung ir/ser & H.P., A.S. & Kontrastierung: ir = Bewegung, ser = Sein \\
\hline
Modalverb-Auswahl (poder/querer/tener que) & A.A., H.P. & Bedeutungsunterschiede visualisieren \\
\hline
Aussprache [b]/[\ss] & Alle Gruppe A & Chorisches Sprechen, Einzelkorrektur \\
\hline
Leseverständnis komplexer Texte & Schwächere Gruppe B & Wörterbuch-Hilfe, Vorentlastung \\
\hline
Unsicherheit bei Assessment-Ankündigung & A.A., A.S., H.P. & Beruhigen, Übungszeit betonen \\
\hline
Zeitproblem (zu viel Stoff) & -- & Aufg. 6 (Phonetik) als HA-Option \\
\hline
\end{tabular}

% ============================================================================
% 9. HAUSAUFGABE
% ============================================================================
\section{Hausaufgabe}

\textbf{Gruppe A:}
\begin{itemize}
    \item Vokabeln S. 32--33 lernen (ir, poder, querer, tener que)
    \item Optional: Aufg. 6 (Aussprache) zu Hause laut üben
    \item \textbf{Vorbereitung Assessment:} Stadtviertel-Beschreibung schriftlich vorbereiten
\end{itemize}

\textbf{Gruppe B:}
\begin{itemize}
    \item Aufg. 3a fertigstellen (falls nicht abgeschlossen)
    \item Vokabeln S. 105--106 wiederholen
    \item \textbf{Vorbereitung Assessment:} Stichpunkte für Vortrag ``Mi profesión ideal'' erstellen
\end{itemize}

% ============================================================================
% 10. ANLAGEN
% ============================================================================
\section{Anlagen}

\begin{enumerate}
    \item Bewertungsbogen Gruppe A: Partnerprüfung ``Mi barrio''
    \item Bewertungsbogen Gruppe B: Kurzvortrag ``Mi profesión ideal''
    \item Konjugationstabelle \textit{ir} (Tafelbild)
    \item Übersicht Modalverben (poder/querer/tener que)
    \item Wortschatzliste Stadtviertel (für Gruppe A)
    \item Lehrwerkseiten S. 32--33, S. 105--106
\end{enumerate}

% ============================================================================
% 11. DIDAKTIK-SCORE
% ============================================================================
\newpage
\section{Didaktik-Score}

\begin{scorebox}
\textbf{Bewertung der didaktischen Qualität nach 10 Dimensionen}\\
\textbf{Ziel-Score: $\geq$ 9.0 (Premium-Qualität)}

\vspace{1em}
\begin{tabular}{|c|l|c|c|p{4.5cm}|}
\hline
\textbf{\#} & \textbf{Dimension} & \textbf{Gew.} & \textbf{Score} & \textbf{Notiz} \\
\hline
1 & Differenzierung & 15\% & 9/10 & Klare A/B-Trennung, niveaugerechte Aufgaben \\
\hline
2 & Taxonomie (AFB) & 10\% & 9/10 & AFB I--III abgedeckt (FZ1--8) \\
\hline
3 & Lernziel-Alignment & 10\% & 9/10 & Passt zu Rahmenplan, Assessment-orientiert \\
\hline
4 & Aktivierung & 10\% & 8/10 & Überwiegend rezeptiv, aber Paararbeit \\
\hline
5 & Scaffolding & 10\% & 9/10 & Tabellen, Beispiele, Lehrerhilfe \\
\hline
6 & Feedback-Qualität & 10\% & 8/10 & Zwischensicherung, aber wenig individuell \\
\hline
7 & Sprachliche Klarheit & 10\% & 9/10 & Klare Anweisungen, Gruppenaufteilung \\
\hline
8 & Motivation \& Relevanz & 10\% & 9/10 & Assessment als Motivator, Alltagsthemen \\
\hline
9 & Inklusion (UDL) & 10\% & 8/10 & Verschiedene Zugänge (Hören, Lesen, Schreiben) \\
\hline
10 & Praktikabilität & 5\% & 10/10 & Realistischer Zeitplan, klare Struktur \\
\hline
\multicolumn{3}{|r|}{\textbf{GESAMT:}} & \textbf{8.8/10} & Knapp unter Ziel \\
\hline
\end{tabular}

\vspace{1em}
\textbf{Bewertung: } \textcolor{warningorange}{8.8/10 -- Iteration empfohlen}

\vspace{1em}
\textbf{Stärken:}
\begin{itemize}
    \item Klare Differenzierung durch Gruppenstruktur (A vs. B)
    \item Assessment-Ankündigung als strukturgebendes Element
    \item Bewertungskriterien frühzeitig transparent
    \item Zeitplan realistisch und mit Puffer
\end{itemize}

\textbf{Verbesserungspotential:}
\begin{itemize}
    \item[$\Box$] \textbf{Aktivierung:} Mehr Sprechzeit für SuS einplanen (Mini-Dialoge vor Sicherung)
    \item[$\Box$] \textbf{Feedback:} Individuelle Rückmeldung während Übungsphasen verstärken
    \item[$\Box$] \textbf{Inklusion:} Auditive Unterstützung für schwächere SuS (Audio-Track nutzen?)
\end{itemize}

\vspace{1em}
\textbf{Besondere Stärke dieser Stunde:} Die \textbf{Assessment-Ankündigung} schafft Transparenz und gibt den SuS Orientierung. Durch das frühzeitige Verteilen der Bewertungskriterien wissen alle SuS, worauf es ankommt -- dies reduziert Prüfungsangst und ermöglicht gezieltes Üben.
\end{scorebox}

% ============================================================================
% 12. REFLEXIONSFRAGEN (nach der Stunde)
% ============================================================================
\section{Reflexionsfragen (nach der Stunde)}

\begin{enumerate}
    \item Haben alle SuS die Assessment-Kriterien verstanden?
    \item Wie war die Reaktion der schwächeren SuS (A.A., A.S., H.P.) auf die Ankündigung?
    \item War die Zeitaufteilung zwischen Gruppe A und B angemessen?
    \item Konnten die ir-Formen ausreichend geübt werden?
    \item Haben die Gruppe-B-SuS die Texte verstanden?
    \item Welche Fragen kamen zur Assessment-Vorbereitung?
\end{enumerate}

\end{document}
